% Chapter 35: Treating Light as a Wave
\chapter{Treating Light as a Wave}

\step{A Counterintuitive Opening}

\begin{mainline}
Blow a soap bubble and watch. The film is transparent and sunlight colorless, yet vivid bands roll over its surface: red, green, blue, and violet flowing like liquids. No one painted or dyed them. A breeze immediately changes the pattern. Where are the colors hiding?

Now tilt an old compact disc toward a lamp. A complete rainbow bursts across it, though its surface is shiny plastic without a drop of pigment. An equally shiny glass plate or mirror cannot produce this rainbow. What does the disc have that the mirror lacks?

Those two observations are only a warm-up. The third is truly unsettling. Send a narrow light beam through two closely spaced slits in a barrier onto a wall. Common sense predicts two bright lines, dark between and on either side, nothing more. In reality, a whole row of evenly spaced bright and dark stripes spreads out like a zebra, a dozen or more. Stranger still: locate the exact center of a \emph{dark} stripe, where there is no light, then block one slit. Less light enters and the wall's total brightness falls, yet that formerly black spot \emph{lights up}.

Block some light and a dark place becomes bright. Every concept in this chapter was forced into existence by that counterintuitive sentence.
\end{mainline}

\step{Make a Guess First}

As always, put your intuition on the table before peeking at the answers.

Question one. At night, look at a distant streetlight, or put your phone's flashlight at maximum brightness across the room. Hold two fingers together with a very narrow gap and look through it. What appears? (a) A squeezed, narrow light; (b) several bright stripes along the fingers; (c) alternating bright and dark stripes arranged \emph{perpendicular to} the fingers. Notice option (c)'s direction. If it is not your intuitive choice, remember the discomfort.

Question two. A double-slit experiment produces a completely dark stripe. Block one slit, halving the incoming light. Does that dark position (a) become darker, (b) remain black, or (c) brighten?

Question three. Does a CD's rainbow arise from some substance in its plastic that changes color with angle? If you polish away the rainbow-producing structure to leave a completely smooth disc, does the rainbow remain?

Write your three answers and reasons. We will settle them individually at the end.

\step{Hands-on Experiment}

\begin{experiment}[Finger-Gap Stripes and a CD Rainbow]
\textbf{Materials}: a phone with a flashlight, an old CD or DVD, and a room you can darken. Optional: dish detergent, water, and wire bent into a loop.

\textbf{Step one}: switch off the room lights and place the phone flashlight two or three meters away, facing you. Bring your index and middle fingers together, squeezing slightly to leave a hair-thin gap. Hold it before your eye and look through at the ``lamp.'' Squint and inspect carefully: instead of one narrow light, a small row of alternating bright and dark stripes extends \emph{perpendicular to the gap}. Squeeze the gap narrower and the stripes spread farther apart. Ray intuition says ``a narrower gap should narrow the view,'' but experiment insists on the opposite.

\textbf{Step two}: shine the flashlight obliquely onto the disc's shiny side, opposite its label, and rotate the disc while watching the reflection. At one angle a rainbow spreads across it; change angle and its arrangement changes. Repeat with a mirror or glass plate: no colors.

\textbf{Step three, optional}: dilute detergent in water, lift a soap film with the wire loop, and hold it vertically while observing reflected lamplight. After a few minutes, colors flow into horizontal bands. If you wait long enough, a \emph{black} region appears at the top. The film has not broken there; this often happens before it bursts.

\textbf{Record}: which direction do the finger-gap stripes extend? Does narrowing the gap crowd or separate them? How does a disc's rainbow differ from a mirror's reflection? Why does the film's top blacken first? Write down your observations; this chapter unpacks them all.
\end{experiment}

Another free observation: a thin oil film floating on a roadside puddle after rain displays brilliant colors. Transparent water, transparent oil, colorless sunlight, yet colors appear. It shares the soap bubble's secret.

\step{The Old Model Runs into Trouble}

Chapters 33 and 34 treated light as rays: straight propagation, reflection from mirrors, refraction at boundaries, and lenses focusing sharp images. This model has an impressive record: cameras, telescopes, spectacles, and optical fibers were all designed with it. Why invite trouble when it works so well?

Because trouble arrives uninvited, and in a whole procession.

First trouble: shadows refuse neat edges. Rays predict knife-sharp boundaries, simply bright or dark. Yet inspect any sunlit shadow and a blurred transition appears. In a pinhole camera, anticipated at Chapter 34's end, shrinking the hole too far makes the image blur rather than sharpen. In the ray model, a smaller aperture should only reduce admitted light, with no reason to change the pattern's shape. But the shape does change.

Second trouble: the direction is wrong. A vertical finger gap should, according to rays, squeeze the visible light vertically; the observed stripes instead spread \emph{horizontally}. Narrowing the gap makes the pattern spread perpendicularly. Straight-line propagation cannot explain light entering ``forbidden directions.''

Third trouble: there is nowhere to put color. A reflecting surface sends rays in one definite direction. A CD is merely many reflecting surfaces, so why should it split white light into a rainbow that changes with angle? Pigment cannot explain it: rotating the disc relocates the rainbow, while material composition does not change with angle.

Fourth and fatal trouble: the double slit. For two beams meeting at the same wall position, the ray model knows only one arithmetic operation: add their brightnesses. The result must brighten or stay unchanged, never darken. Yet experiment produces stable, completely dark stripes, which brighten when one beam is blocked. ``Light plus light equals darkness'' is unwriteable in the ray model's language, as jarring as \(1+1=0\).

These four problems share a feature: the relevant obstacles, finger gaps, pinholes, and disc grooves, are all very \emph{small}. Rays triumph on large scales and retreat on small ones. This is not a minor numerical error; the model omits a fundamental aspect of light. To restore it, return to Chapter 17's old friend: waves.

\step{Inventing New Concepts}

\begin{mainline}
\textbf{First concept: light is a wave with a wavelength.} Chapter 25 established that light is an electromagnetic wave traveling through vacuum at speed \(c\). Waves have two basic identifiers: frequency, oscillations per second, and wavelength, the spatial length of one complete cycle. Visible wavelengths run roughly from \(400\) nanometers for violet to \(700\) nanometers for red. A hair's diameter spans a little over a hundred visible wavelengths. Colors are wavelength's roster: longer looks red, shorter violet. We normally miss light's ``waviness'' because its wavelength is so short that it behaves like a straight line on everyday scales. Only when obstacles approach a wavelength, or exceed it by no more than a few hundred times, can its wave nature no longer hide.

\textbf{Second concept: Huygens' principle, every wavefront point is a new source.} Wherever a wave has reached, each point of its wavefront, a surface joining points at the same phase at one instant, can be treated as a tiny source emitting forward spherical wavelets. The next wavefront is the envelope formed by their superposition. What does this resemble? A stadium wave: nobody runs around the entire stadium; each spectator merely rises and sits in response to a neighbor, yet the wave pattern circles the crowd. What does it not resemble? An actual row of little lamps scattered across the front. Wavelets are not independent objects but a way to calculate where the wave goes next.

Before explaining oddities, the principle must explain the ordinary: why do waves usually go straight? Imagine a broad plane wavefront emitting spherical wavelets at every point. Sideways wavelets arrive from dense neighboring points with mismatched phases and cancel, while directly ahead they arrive together and reinforce a new plane front. The envelope remains flat, so the wave advances straight. On everyday scales, it is this wavelet envelope that travels straight; Chapter 33's ``ray'' is merely a reference line perpendicular to it. Reflection and refraction also follow from wavelet geometry, though we will not derive them here. The point is that Huygens' principle does not overthrow rays but explains them as the broad-wavefront limit of waves.

The principle's real power emerges at small scales. It explains bending in one sentence: when a wavefront passes an obstacle's edge, nearby unblocked points continue emitting spherical wavelets, some of which naturally enter the geometrical shadow. The smaller the obstacle, the narrower the admitted wavefront; that surviving portion increasingly resembles a point source spreading spherical waves in every direction. This is diffraction. Light does not ``bounce sideways'' at the edge; the straight-line approximation has expired at small scales.

\textbf{Third concept: phase and superposition, adding waves by their rhythm.} A wave needs not only ``how strongly it oscillates,'' amplitude, but ``where it is in its cycle,'' phase. Where two waves meet, the local oscillation is their sum: crest meets crest and matching phases double the vibration; crest meets trough and opposite phases cancel. Chapter 18 let you hear this: between two speakers playing the same song are unusually quiet spots. Different path lengths bring their sound waves there in opposite phase. Transfer the idea to light. Light also oscillates, but its electric and magnetic fields oscillate, not a material displacement. Unlike sound and water waves, it needs no carrying ``surface.'' Two light waves therefore add and subtract by phase too: matching phases produce bright stripes, opposite phases dark stripes. ``Light plus light equals darkness,'' impossible for rays, becomes half a sentence in wave language: crest plus trough.
\end{mainline}

\begin{center}
\begin{tikzpicture}[scale=1.0]
  % Barrier and double slit
  \draw[very thick] (-2.6,-1.8) -- (-2.6,-0.35);
  \draw[very thick] (-2.6,-0.15) -- (-2.6,0.15);
  \draw[very thick] (-2.6,0.35) -- (-2.6,1.8);
  \node[left] at (-2.6,-0.25) {Slit};
  % Incident wavefronts
  \foreach \x in {-4.6,-4.0,-3.4} {\draw[blue!60!black] (\x,-1.5) -- (\x,1.5);}
  \draw[-{Stealth[length=2.5mm]},blue!60!black] (-4.6,0) -- (-3.0,0);
  \node[blue!60!black,above] at (-3.9,1.5) {Incident plane wave};
  % Wavelets from the slits (arcs)
  \foreach \r in {0.5,1.0,1.5,2.0} {
    \draw[gray!60,dashed] (-2.6,0.25) ++(0:\r) arc (0:55:\r);
    \draw[gray!60,dashed] (-2.6,0.25) ++(0:\r) arc (0:-55:\r);
    \draw[gray!60,dashed] (-2.6,-0.25) ++(0:\r) arc (0:55:\r);
    \draw[gray!60,dashed] (-2.6,-0.25) ++(0:\r) arc (0:-55:\r);}
  % Screen and fringes
  \draw[very thick] (2.6,-1.8) -- (2.6,1.8);
  \foreach \y/\w in {-1.4/0.5,-0.9/0.9,-0.45/0.5,0/0.9,0.45/0.5,0.9/0.9,1.4/0.5}
    {\fill[red!60!black] (2.66,\y-0.09) rectangle (2.78,\y+0.09);}
  \node[below] at (2.7,-1.8) {Bright screen fringes};
  % Two paths to one screen point
  \draw[red!70!black,thick] (-2.6,0.25) -- (2.6,0.9);
  \draw[red!70!black,thick] (-2.6,-0.25) -- (2.6,0.9);
  \node[red!70!black,right] at (0.4,0.75) {Two unequal paths};
\end{tikzpicture}
\end{center}

\begin{mainline}
\textbf{Fourth concept: optical path difference judges brightness.} The two slits are illuminated by the same wavefront and serve as in-phase wavelet sources. A point on the wall is generally at unequal distances from them. Waves leave in phase and travel different paths, so their phase difference on arrival is entirely determined by path difference. An integer number of wavelengths brings crest to crest: bright. Half, one-and-a-half, two-and-a-half wavelengths, and so on bring crest to trough: dark. The zebra stripes are simply a map of path differences. This also explains why the slits must be close and share one illuminating beam: the difference must be only a few wavelengths for phase differences to remain stable and fringes visible.

\textbf{Fifth concept: coherence, not any two beams produce stable interference.} Why can two flashlights never cast dark interference stripes on a wall? Cancellation must persist long enough to be seen. Ordinary sources, Sun, bulbs, and flashlights, emit unrelated short wave trains. The phase relation between two beams randomly changes ten quadrillion times per second, making bright and dark fringe positions flicker wildly. Eyes and cameras capture only the blurred average, uniform brightness. Sources with a stable phase relationship are \emph{coherent}. The double-slit trick divides the \emph{same} wavefront into two routes. However the train's phase randomly changes, both routes change identically, leaving their difference fixed and the stripes steady. Laser pointers enable home interference demonstrations because their wave trains are much more orderly and their coherence times much longer. Two flashlights' missing stripes do not reflect inadequate instruments; those beams lack a stable mutual rhythm.
\end{mainline}

\step{The Model in One Sentence}

\begin{oneline}
Treat light as a wave: brightness at each screen point depends on whether oscillations arriving by different paths are in or out of step. Whole-wavelength path differences reinforce; half-wavelength differences cancel. As obstacle size approaches wavelength, light ceases traveling straight and spreads according to Huygens' principle.
\end{oneline}

\step{The Mathematical Translator}

\begin{mainline}
First express double-slit fringes mathematically. Let slit separation be \(d\), the distant screen be \(L\) away, and the observed point lie at angle \(\theta\) from the central direction. Geometry gives the path difference
\[
  \delta = d\sin\theta .
\]
The judging rule is already known: whole wavelengths give brightness, half-offset wavelengths darkness. Hence
\[
  d\sin\theta = m\lambda \ \text{(bright)},\qquad
  d\sin\theta = \left(m+\tfrac12\right)\lambda \ \text{(dark)},\qquad m=0,1,2,\dots
\]
Read the terms: the left is ``how far the paths differ,'' the right ``how many wavelengths.'' The direction of a fringe is \(\theta\), and light's own ruler is \(\lambda\). At small angles, neighboring bright fringes are separated on the screen by
\[
  \Delta y \approx \frac{L\lambda}{d}.
\]
Three routine checks follow. First, at \(\theta=0\), the center, \(\delta=0\), satisfying the bright condition for every wavelength. White illumination therefore always gives a white central fringe: every color is in phase there. Second, decrease \(d\): \(\Delta y\) increases and stripes spread. As \(d\) tends to zero, their spacing tends to infinity and the whole screen becomes uniformly bright. The slits merge into one source and interference disappears: sensible. Third, let \(\lambda\) tend to zero relative to \(d\): \(\Delta y\) vanishes, the eye cannot separate the dense fringes, and only the two slits' geometrical images remain. Waves automatically reduce to Chapter 33's rays. Incorporating the old model as a limit is the new model's first certificate of competence.

% BEGIN TRANSLATOR SAFETY NOTE
\textbf{Translator safety note:} Do not look into the beam or its reflections; laser pointers are not toys. Use a supervised setup that keeps all beams away from people and animals, or treat this as a calculation. See \href{https://www.fda.gov/consumers/consumer-updates/laser-toys-how-keep-kids-safe}{FDA laser safety guidance}.
% END TRANSLATOR SAFETY NOTE

Try real numbers. A red laser pointer has wavelength about \(650\) nanometers. Pierce two holes about \(0.25\) millimeters apart in opaque paper, with a screen \(1\) meter away:
\[
  \Delta y \approx \frac{1\ \text{m}\times 650\times10^{-9}\ \text{m}}{0.25\times10^{-3}\ \text{m}}
  \approx 2.6\ \text{mm}.
\]
Millimeter-scale fringes are just large enough to count by eye, so this works in a dark living room. Conversely, a widely spaced pair with \(d=5\) millimeters gives only about \(0.13\) millimeters between fringes, blurring them together. Interference fringes are always present; usually they are simply too fine to see.

The second equation describes single-slit diffraction. Huygens' principle treats a slit of width \(a\) as a row of wavelet sources. Compare the ends: when their path difference reaches one wavelength, corresponding pairs within the slit cancel, producing the first dark fringe:
\[
  a\sin\theta = \lambda \quad\text{(first dark fringe)},
\]
with central bright-fringe angular width about \(2\lambda/a\). Check: larger \(a\) gives smaller \(\theta\), so a broad slit barely diffracts and light goes straight, recovering rays. When \(a=\lambda\), \(\sin\theta=1\) and \(\theta=90^\circ\): the first minimum moves to the horizon and light spreads uniformly throughout the forward region. An aperture as small as a wavelength becomes a spherical source. Between those cases, smaller \(a\) means a broader pattern. This completely explains ``squeeze tighter, spread farther'' and why the pattern extends \emph{perpendicular to the slit}: diffraction occurs only along the dimension being narrowed.
\end{mainline}

\begin{center}
\begin{tikzpicture}[scale=1.0]
  % Single-slit diffraction intensity
  \draw[-{Stealth[length=2.5mm]}] (-4.2,0) -- (4.4,0) node[below right] {\shortstack{Screen\\position}};
  \draw[-{Stealth[length=2.5mm]}] (0,0) -- (0,2.4) node[below left] {Brightness};
  \draw[very thick,blue!70!black,smooth] plot coordinates {
    (-4.0,0.00) (-3.6,0.03) (-3.2,0.00) (-2.9,0.06) (-2.6,0.00)
    (-2.2,0.15) (-1.9,0.00) (-1.4,0.50) (-0.8,1.40) (0,2.00)
    (0.8,1.40) (1.4,0.50) (1.9,0.00) (2.2,0.15) (2.6,0.00)
    (2.9,0.06) (3.2,0.00) (3.6,0.03) (4.0,0.00)};
  \draw[dashed] (0,0) -- (0,2.0);
  \draw[dashed] (1.9,0) -- (1.9,0.25);
  \draw[dashed] (-1.9,0) -- (-1.9,0.25);
  \node[below] at (1.9,0) {First minimum};
  \node[below] at (-1.9,0) {First minimum};
  \draw[{Stealth[length=2mm]}-{Stealth[length=2mm]},red!70!black] (-1.9,1.0) -- (1.9,1.0);
  \node[red!70!black,above] at (0,1.0) {Central fringe width \(\approx 2\lambda/a\)};
\end{tikzpicture}
\end{center}

\begin{mainline}
% BEGIN TRANSLATOR SAFETY NOTE
\textbf{Translator safety note:} Keep this laser calculation theoretical unless a competent supervisor provides an eye-safe demonstration setup. Do not view the beam directly or through the hair. See \href{https://www.fda.gov/consumers/consumer-updates/laser-toys-how-keep-kids-safe}{FDA laser safety guidance}.
% END TRANSLATOR SAFETY NOTE

Estimate again using a hair. Place one of diameter about \(70\) micrometers before a laser pointer, with a screen \(1\) meter behind. A bright line lies in its shadow, flanked by diffraction fringes. The first minimum occurs at angle approximately
\[
  \theta \approx \frac{\lambda}{a} = \frac{550\times10^{-9}\ \text{m}}{70\times10^{-6}\ \text{m}}
  \approx 0.008\ \text{radians},
\]
corresponding on a screen \(1\) meter away to about \(8\) millimeters from the center. A hair becomes a wavelength-measuring instrument: measure fringe positions and infer a wavelength of a few hundred nanometers.

The third equation describes a grating. Adjacent tracks in a CD's spiral groove are about \(1.6\) micrometers apart, a ruler bearing thousands of ``slits,'' called a diffraction grating. Many equally spaced slits sharpen interference fringes enormously, with bright directions satisfying
\[
  d\sin\theta = m\lambda ,
\]
exactly the double-slit form, except that the bright fringes become narrower and brighter. Insert real values: \(d=1.6\) micrometers. In the first order, \(m=1\), violet light at \(400\) nanometers gives \(\theta\approx 14^\circ\), while red at \(650\) nanometers gives \(\theta\approx 24^\circ\). Within one order, different colors go in different directions, separating white light into a rainbow. A mirror lacks grooves and this ruler, so it produces no colors. Was your third guess right? The rainbow lives not in the plastic but in the track spacing. Smooth away the grooves and the rainbow dies.

The key ``path difference determines brightness'' opens more than optics. Noise-canceling headphones industrialize it in acoustics: a microphone hears external noise, a chip generates an opposite-phase sound for your ear, and the two cancel at the eardrum. ``Sound minus sound equals quiet'' and ``light minus light equals dark'' are the same principle. The difference is light's much shorter wavelength: optical ``noise-canceling headphones'' would require controlling two optical paths to within one wavelength. That is why interferometers, including laser instruments detecting gravitational waves, are difficult to build.
\end{mainline}

\begin{mathbridge}
Phase expresses ``in step adds, out of step cancels'' more cleanly. Path difference \(\delta\) corresponds to phase difference
\[
  \Delta\varphi = \frac{2\pi}{\lambda}\,\delta ,
\]
so every extra wavelength adds one full turn, \(2\pi\). Superposing equal-amplitude waves gives a resultant amplitude proportional to \(\cos(\Delta\varphi/2)\), while brightness, or intensity, is proportional to amplitude squared:
\[
  I = I_0\cos^2\!\frac{\Delta\varphi}{2}.
\]
Check: \(\Delta\varphi=0\) gives \(I=I_0\), the brightest case; \(\Delta\varphi=\pi\) gives \(\cos(\pi/2)=0\), complete darkness. Between lies a smooth transition, explaining graded rather than knife-edged fringes. Energy conservation hides here too: average \(I\) over the whole screen and it equals the sum of the separate beams' intensities. Every penny ``missing'' from dark places is piled into bright ones.

Thin-film interference uses the same key, replacing geometric distance with optical path length. In a film of index \(n\), light crosses thickness \(t\). Since wavelength inside shrinks to \(\lambda/n\), phase accumulates according to \(nt\), and reflections from the two surfaces differ in optical path by \(2nt\). One detail remains: reflection from an optically rarer into a denser medium adds a half-wavelength phase reversal, a half-wave loss. Only the soap film's upper reflection has it. Destructive reflection therefore satisfies
\[
  2nt = m\lambda .
\]
Antireflection coatings on camera and spectacle lenses follow this equation, choosing thickness for cancellation near \(2nt=\lambda/2\), or \(t\approx \lambda/(4n)\). Magnesium fluoride of index about \(1.38\), designed for green \(550\)-nanometer light, needs a film about \(100\) nanometers thick, roughly a virus's size. This nanoscale layer reduces reflection, increases transmission, and suppresses stray light. A final challenge will let you explain its purple appearance in white light.
\end{mathbridge}

\begin{challenge}
The recurring \(\lambda/a\) is one corner of a larger map. In Fourier language, the more tightly an aperture, slit, hair, pupil, or telescope opening, confines light, the broader its angular spread. Spatial compression trades for directional spreading; the two are Fourier-transform partners. Narrower slits producing broader diffraction and shorter sounds producing less certain frequencies are two faces of one mathematical theorem. In the quantum section it returns in grander form: photons spreading in direction after a narrow slit anticipate the position--momentum uncertainty relation in Chapter 43. It reflects wave nature, not crude instruments.

Fourier optics offers another elegant intuition: \emph{a lens turns direction into position}. Parallel light incident at one angle focuses at one focal-plane point; change direction and that point shifts. A lens faithfully translates ``which way the diffraction spreads'' into ``where it lands in the focal plane,'' drawing the aperture's Fourier transform there in real time. Far-field diffraction, spectrum analysis, and Chapter 37's imaging theory grow from this one sentence.

This map also condemns every optical instrument to a limit: any circular aperture of diameter \(D\) spreads a point source into a bright spot of angular radius about \(1.22\lambda/D\). Two sources closer than that cannot be resolved, the Rayleigh criterion. For a human pupil of about \(3\) millimeters and green light of \(550\) nanometers,
\[
  \theta_{\min} \approx \frac{1.22\times 550\times10^{-9}\ \text{m}}{3\times10^{-3}\ \text{m}}
  \approx 2\times10^{-4}\ \text{radians},
\]
about one arcminute, matching the \(1.0\) line on an eye chart. Your resolution limit is not insufficient retinal precision but diffraction at the pupil. Likewise, chip lithography uses extreme ultraviolet of wavelength \(13.5\) nanometers: finer circuits require a shorter ``light ruler.'' This is diffraction's multibillion-dollar face in the semiconductor industry.
\end{challenge}

\step{The Model's Limits}

\begin{mainline}
Waves have won every battle here, but mark their borders before crossing them elsewhere.

Limit one: \emph{this is still not light's final face}. Reduce intensity until, on average, only one ``small portion'' reaches the detector at a time. Waves predict energy continuously and evenly washing over the screen, but the detector registers discrete clicks. In a weak-light double-slit experiment, individual bright dots accumulate before gradually revealing the fringes. Waves correctly predict fringe \emph{positions} but cannot explain arrival one grain at a time. This is not an instrument defect: light has a third face, belonging to the quantum section from Chapter 42 onward. Chapter 37 first clarifies the relation among all three.

Limit two: \emph{wave optics matters where wavelength cannot be ignored}. For obstacles thousands or tens of thousands of wavelengths across, diffraction angle \(\lambda/a\) is too small to measure and waves and rays predict the same result. Use rays then; they are simpler. Light has not ``turned back into particles''; the appropriate approximation has arrived for duty.

Limit three: \emph{coherence governs visible interference}. Ordinary thermal sources often have coherence lengths of only micrometers; slightly larger path differences wash fringes out. Sunlit soap films are thin enough that optical differences remain within this length, so colors appear. The same sunlight performs poorly in a double-slit experiment because slit-to-screen path differences exceed the wave train's ``memory.''
\end{mainline}

Three characteristic misuses also deserve listing:

\begin{itemize}[leftmargin=1.8em]
  \item \emph{``Energy vanishes at dark fringes.''} Wrong. Interference reallocates energy: bright places gain what dark places lose, and total screen energy equals the two beams' sum. Cancellation never violates conservation, because it cancels \emph{oscillations}, not energy.
  \item \emph{``Diffraction is light bouncing or scraping sideways at aperture edges.''} Wrong. Huygens' principle contains no collisions: the surviving wavefront keeps emitting wavelets that naturally spread forward. Enlarge the aperture tenfold and its edge remains, yet diffraction disappears. The key is aperture-to-wavelength ratio, not the edge itself.
  \item \emph{Mistaking a wave graph for light's flight path.} The textbook's undulating line represents electric-field strength and direction at different positions, not a wavy route through space. Light still travels straight in a uniform medium; the electromagnetic field oscillates, not the flight path. Likewise, double-slit interference is not ``two beams colliding and fighting,'' but their local phase-sensitive sum at each screen point.
\end{itemize}

\step{Explaining the Opening Phenomenon}

Now collect the outstanding accounts with the new model.

\textbf{Soap-bubble colors.} A soap film is a thin water layer. Some incident sunlight reflects at its upper surface; another portion enters, reflects below, and emerges again. These beams follow different routes with optical path difference about \(2nt\). For a wavelength satisfying constructive interference, that color is reinforced in reflection while others are weakened. Gravity makes thickness vary from top to bottom, arranging favored colors by height into horizontal bands. Wind perturbs thickness and makes the patterns flow. White light contains every wavelength, so each position can find its own color. Most striking is the \emph{black} region at the top. When thickness becomes far smaller than wavelength, \(2nt\) approaches zero, but the upper reflection has a half-wave phase reversal while the lower does not: the reflections are naturally opposite. At extreme thinness they cancel, making the film black in reflection. Transparent water becomes black through ``light minus light.'' Blackness is not dirt but interference's signature.

Roadside oil films use the same principle, with uneven thickness making more irregular colors. Nature has mastered this trick: peacock blue-green plumage, metallic flashes on dragonfly wings, and some beetle-shell colors arise not from pigment but interference in nanoscale layered structures, called structural color. It has a talent pigment can never learn: changing color with viewing angle because the optical path difference changes.

\textbf{The disc's rainbow.} Already settled by the grating equation, \(d\sin\theta=m\lambda\): \(1.6\)-micrometer track spacing sends colors in different directions. No grooves, no rainbow; polish them away and it disappears.

\textbf{Finger-gap stripes.} The fingers form a narrow slit. Single-slit diffraction spreads lamplight along its narrowed dimension into bright and dark stripes, with angular width about \(2\lambda/a\). Narrower means wider spread. Question one's answer is (c), opposite the intuitive direction. Tiny gaps between eyelashes while squinting at a streetlight produce similar rays.

\textbf{A double-slit dark spot brightens.} At a dark fringe center, opposite-phase oscillations from the two paths cancel exactly. Block one slit and the remaining wave dances alone, with nothing to cancel it. The point receives the uniform brightness of single-slit diffraction. Question two is also (c): less light, but a brighter point. Total screen energy does decrease, yet redistribution gives the previously dark position a share. The counterintuitive part is not the experiment but the old assumption ``more light makes every place brighter.''

\step{Thought Experiments and Challenges}

\begin{thought}[Make the Whole Earth a Telescope]
Rayleigh's criterion says finer detail requires a shorter wavelength or a larger aperture \(D\). Astronomers pushed this to an audacious extreme: since an Earth-sized lens is impossible, telescopes scattered worldwide observe one target \emph{simultaneously}, precisely recording arriving waves. Later, align their phases and combine them interferometrically, an engineering version of this chapter's phase-sensitive addition. In 2019, the Event Horizon Telescope used a global radio array with baselines exceeding ten thousand kilometers at wavelength \(1.3\) millimeters. Its effective aperture approached Earth's diameter and resolution reached about twenty microarcseconds, like resolving a coin in Shanghai from Beijing, finally imaging the shadow of the black hole in galaxy M87. That famous orange ring is fundamentally a diffraction-limited image developed by interference with Earth as the aperture. From soap bubbles to black holes, only ``optical path difference'' lies between.
\end{thought}

\begin{thought}[One Photon at a Time]
Dim the double-slit source until, on average, at most one ``portion'' of light travels through the apparatus at once. Waves predict continuously fading fringes. Experiment instead shows one discrete screen point at a time, with unpredictable position. Thousands accumulate into precisely the double-slit pattern, as if every portion knew about both slits and chose its arrival according to the fringe probability map. Stranger still, a detector beside the slits recording which route it took immediately destroys the fringes, leaving two geometrical bright patches. Neither classical waves nor particles suffice. Light does not transform back and forth between them; it requires a new language of quantum-state superposition and probability amplitudes. This classical-physics disaster erupts formally in Chapter 42 and is retried with quantum mechanics from Chapter 43. For now remember: the fringe pattern's phase accounts apply even to light arriving portion by portion. That is their deepest feature.
\end{thought}

\begin{puzzles}
\item At a double-slit dark fringe, two waves cancel. Someone says: ``Cancellation destroys some of their energy; otherwise how could bright fringes exceed the simple sum of the beams?'' Identify the mistake. \emph{Hint}: average the graded fringe intensity over the screen and compare it with the sum of separate intensities. Cancellation affects oscillation at \emph{that point}; energy moves from dark places to bright ones, with the total unchanged.
\item With white sunlight illuminating a double slit, what color is the central fringe? In the first colored band outward, is red or violet closer to the center? \emph{Hint}: at \(\theta=0\), every wavelength has zero path difference and matching phase. Spacing \(\Delta y=L\lambda/d\) is proportional to wavelength. Which is longer, red or violet?
\item On a building's shaded side, you cannot see a rooftop lamp but can receive FM radio clearly. Both are electromagnetic waves; why such different treatment? \emph{Hint}: diffraction angle is about \(\lambda/a\). Visible wavelengths are under a micrometer, FM wavelengths about three meters, and the building tens of meters across. Compare wavelength-to-obstacle ratios and decide which wavelets can spread through the shadow.
\item A lens coating is designed to cancel reflected green light at \(550\) nanometers. Viewed obliquely under white light, its reflection often looks pale purple. Why purple rather than green? \emph{Hint}: green is suppressed most completely, red and blue-violet only partly. What color remains when green is removed? Also ask why manufacturers choose green: to which color are human eyes most sensitive?
\end{puzzles}

We have now met two of light's three faces: Chapter 33's rays and this chapter's waves. Even the wave face is incomplete. Light is transverse, giving its oscillation another dimension, orientation, which introduces polarization and light's journey inside matter in Chapter 36. Unifying the first two faces and understanding why the third is necessary belong to Chapter 37 and the quantum section. A disc, a soap bubble, and a finger gap have already taken ``what is light?'' much deeper.
